% Physical unit 008; accepted Session 02.
\GordanUnitStart{008}{printed page 3}
\GordanTopFolio{3}
aliquem immobilem maxime tractari solent, sic in sphaeroide curvae arcum \gq{$s$} (quem, quod $\psi$ longitudo est, extensionem nominare placet); et angulum \gq{$T$} (Azimut), quem curva in punctis, in quibus meridianum secat, cum eo efficit, accuratius considerare volumus.

\begin{center}\S\ 9.\end{center}
Placet igitur quantitates $s$ et $T$ referre ad quantitates $l$ et $\psi$. Quare consideremus in aliqua curva sphaeroica duo puncta $A(l,\psi)$ et $B(l+dl,\psi+d\psi)$, quibus nihil propius cogitari potest. Secet meridianus puncti $A$ circulum puncti $B$ in puncto $C(l+dl,\psi)$. Quibus punctis triangulus $ABC$ formatus, cum minimus sit, in plano situs sumi potest. Ejus latera $AC$ et $BC$, quorum alterum meridianum, alterum circulum tangit, inclinata sunt sub angulo recto.

Itaque:
\[
 {BC\over AC}=\tg BAC;\qquad {AC\over AB}=\cos BAC;\qquad {BC\over AB}=\sin BAC.
\]
\[
  AB^2=AC^2+BC^2.
\]
Atqui:
\[
\begin{gathered}
 BAC=T;\qquad AB=ds,\\
 BC\;\hbox{(arcus circuli radio \gq{$a\cos l$} descripti)}=a\cos l\,d\psi,\\
 AC\;\hbox{(arcus ellipseos)}=a\sqrt{1-e^2\cos^2 l}\,dl.
\end{gathered}
\]
Quibus valoribus superiorum loco positis, fit:
\[
 ds^2=a^2(1-e^2\cos^2l)\,dl^2+a^2\cos^2l\,d\psi^2,
\]
\[
\tg T={\cos l\over\sqrt{1-e^2\cos^2l}}{d\psi\over dl};\qquad
\cos T=dl\sqrt{\frac{1-e^2\cos^2l}{(1-e^2\cos^2l)dl^2+\cos^2l\,d\psi^2}},
\]
\[
\sin T=d\psi\sqrt{\frac{\cos^2l}{(1-e^2\cos^2l)dl^2+\cos^2l\,d\psi^2}}.
\]

\begin{center}\S\ 10.\end{center}
In terram spectant e calculis ab ill. Bessel propositis rationes hae:
\[
\begin{array}{rlrl}
 a&=3272077,1399\ \hbox{tois}; & b&=3261139,3284\ \hbox{tois},\\
 \log a&=6,5148235337; & \log b&=6,5133693539,\\
 \log {b\over a}&=9,9985458202; & \log {a\over b}&=0,0014541798,\\
 \varepsilon={a-b\over a+b}&={1-\sqrt{1-e^2}\over 1+\sqrt{1-e^2}}=0,001674184767,\\
 \log e&=8,9122052075; & \log\varepsilon&=7,2238033861.
\end{array}
\]
